\documentclass[11pt]{article}

\usepackage[a4paper,margin=1in]{geometry}
\usepackage{fontspec}
% AI : reference Latin Modern by filename (tectonic bundles it); NoCommon ligatures keeps
% PDF text extraction byte-exact for quote grounding
\setmainfont{lmroman10-regular.otf}[
  BoldFont = lmroman10-bold.otf,
  ItalicFont = lmroman10-italic.otf,
  BoldItalicFont = lmroman10-bolditalic.otf,
  Ligatures = NoCommon]
\usepackage{titlesec}
\usepackage{parskip}

% AI : wording kept identical to the markdown so the assessment step lines up with the rules
\titleformat{\section}{\large\bfseries}{}{0em}{}
\titlespacing*{\section}{0pt}{16pt}{6pt}

\begin{document}

\begin{center}
  {\LARGE\bfseries Company Profile: MediFlow AI}\\[0.6em]
  \rule{\textwidth}{0.4pt}
\end{center}

\noindent\textbf{Sector:} Digital health / clinical decision support\\
\textbf{Headquarters:} Berlin, Germany\\
\textbf{Stage:} Series A

\medskip
MediFlow AI builds an AI-powered symptom triage and care-coordination platform used by
patients and clinics across the European Union and, more recently, the United States. The
following is a summary of the company's current data and operational practices, prepared
for compliance assessment.

\section{Product and Data Collected}
MediFlow's mobile app collects each user's name, email, date of birth, self-reported
symptoms, medication history, and continuous heart-rate data from connected wearables. At
sign-up, users tick a single checkbox agreeing to the Terms of Service and Privacy Policy,
which covers analytics, marketing, product improvement, and health processing in one
combined statement.

Beyond what the triage feature needs today, the app also collects users' GPS location and
phone contact lists ``to support future social and location-based features that are still
on the roadmap.'' This data is currently unused.

\section{Users and Scale}
The platform has roughly 50,000 active users, the majority of whom are residents of the
European Economic Area. Health data is stored indefinitely so that the AI models can be
retrained on the full historical dataset; no automated deletion or retention cut-off is
currently configured.

\section{Infrastructure and Security}
All application traffic is served over HTTPS using TLS 1.3. The primary production database
is hosted in AWS us-east-1 (Northern Virginia) for cost and latency reasons, and European
user records are not segregated from this US deployment. Data stored in the database is not
currently encrypted at rest; the team has flagged this as a backlog item.

Engineers access systems through role-based permissions. Multi-factor authentication is
enforced for administrator accounts but is optional for regular engineers. The company's
most recent third-party penetration test was completed eighteen months ago; a follow-up has
not yet been scheduled.

\section{AI and Automated Decisions}
The triage engine uses a large language model to suggest a likely cause and an urgency level
for each user's symptoms. When the model flags a case as urgent, it is queued for a human
clinician to review before any care recommendation is sent to the user; routine cases
receive an automated explanation of the reasoning. The company keeps a registry of the three
models running in production but has not yet performed a documented bias assessment on any
of them.

\section{Vendors}
MediFlow relies on three sub-processors: an email provider, a cloud analytics tool, and an
LLM inference provider. Signed data processing agreements are in place with all three. The
company does not currently publish a public list of these sub-processors.

\section{Incident Handling and User Requests}
The internal incident-response runbook states that confirmed data breaches are to be
reported to the relevant authority ``within five business days of confirmation.'' Users can
request a copy of their data through support, and the team typically fulfils export and
access requests within 14 days. Deletion requests are honoured within 30 days, although
encrypted database backups containing the deleted records are retained for a further 90 days
before rotation.

\section{Governance}
MediFlow maintains an internal record of its processing activities, reviewed quarterly by
the engineering lead. It has not appointed a dedicated Data Protection Officer; privacy
responsibilities are currently shared informally between the CTO and the head of legal.

\end{document}
